% =============================================================================
%  Robótica (MEEC, IST) — Lab 2: swing de golfe com o UFactory Lite 6
%  Compilar com:  latexmk -pdf main.tex
%  Figuras em Photos/ (cópias de ../*.png; refazer com os scripts de ../../src)
%  Bibliografia em refs.bib (bibtex, estilo IEEEtran)
%  RASCUNHO: procurar \todo{...} — pontos por confirmar antes de entregar.
% =============================================================================
\documentclass[11pt,a4paper]{article}

\usepackage[utf8]{inputenc}
\usepackage[T1]{fontenc}
\usepackage[portuguese]{babel}
\usepackage[margin=2.5cm]{geometry}
\usepackage{graphicx}
\usepackage{booktabs}
\usepackage{tabularx}
\usepackage{float}
\usepackage{amsmath}
\usepackage{siunitx}
\usepackage{caption}
\usepackage[table]{xcolor}
\usepackage{listings}
\usepackage{enumitem}
\usepackage{titlesec}
\usepackage[most]{tcolorbox}
\usepackage[colorlinks=true, linkcolor=azul, citecolor=verde!85!black,
            urlcolor=azul]{hyperref}

\sisetup{output-decimal-marker={,}, per-mode=symbol}

% ----------------------------------------------------------------------------
% Paleta (a mesma do Lab 1): verde para a estrutura, azul para o detalhe, e
% duas cores de sinal usadas só no comentário final (o que correu bem / mal).
% ----------------------------------------------------------------------------
\definecolor{verde}{HTML}{1B7F3B}
\definecolor{azul}{HTML}{1F3A5F}
\definecolor{bom}{HTML}{1B7F3B}
\definecolor{mau}{HTML}{B03A2E}
\colorlet{tabcab}{verde!16}      % linha de cabeçalho das tabelas
\colorlet{tabgrupo}{azul!8}      % linhas que separam grupos dentro das tabelas

\setlist{nosep, leftmargin=*}
\setlist[itemize]{label=\textcolor{verde}{\textbullet}}

\titleformat{\section}{\Large\bfseries\color{verde}}{\thesection}{0.6em}{}
            [{\color{verde}\titlerule[0.8pt]}]
\titleformat{\subsection}{\large\bfseries\color{azul}}{\thesubsection}{0.6em}{}
\titleformat{\paragraph}[runin]{\normalsize\bfseries\color{azul}}{}{0pt}{}

\captionsetup{labelfont={bf,color=azul}, font=small}

\newtcolorbox{caixaresumo}{colback=azul!4, colframe=azul!45, boxrule=0.5pt,
  arc=2pt, left=8pt, right=8pt, top=6pt, bottom=6pt, parbox=false}
\newtcolorbox{destaque}{colback=verde!7, colframe=verde, boxrule=0pt,
  leftrule=3pt, arc=0pt, left=8pt, right=8pt, top=5pt, bottom=5pt}

\renewcommand{\topfraction}{0.9}
\renewcommand{\textfraction}{0.07}
\renewcommand{\floatpagefraction}{0.8}

\lstset{basicstyle=\ttfamily\footnotesize, breaklines=true,
        frame=single, framesep=4pt, rulecolor=\color{azul!45},
        backgroundcolor=\color{azul!4}, keywordstyle=\color{azul}\bfseries,
        stringstyle=\color{verde!85!black}, commentstyle=\color{black!55},
        columns=flexible, keepspaces=true, showstringspaces=false,
        literate={ç}{{\c c}}1 {ã}{{\~a}}1 {é}{{\'e}}1 {ó}{{\'o}}1 {í}{{\'i}}1}
\renewcommand{\lstlistingname}{List}

% \fig{ficheiro}{largura}: inclui a figura se existir; senão deixa um aviso
\newcommand{\fig}[2]{\IfFileExists{#1}{\includegraphics[width=#2]{#1}}%
  {\fbox{\parbox{#2}{\centering\ttfamily\small figura em falta: \detokenize{#1}}}}}
% \todo{...}: ponto do rascunho ainda por confirmar (apagar antes de entregar)
\newcommand{\todo}[1]{\textcolor{mau}{\textbf{[A confirmar: #1]}}}

\newcolumntype{L}{>{\raggedright\arraybackslash}X}
\newcommand{\T}[2]{{}^{#1}T_{#2}}

% Secção "Um segundo swing, desenhado nas juntas" (V3). Só entra no relatório se
% for repetido no robô com 30 mm de folga à mesa (run_robot_v3.py --bola-demo
% --dj2 -6). Para a incluir: trocar \segundoswingfalse por \segundoswingtrue.
\newif\ifsegundoswing
\segundoswingfalse

\title{{\color{verde}\textbf{Golf swing with the UFactory Lite 6}}\\[2mm]
       {\large\color{azul} Robótica --- Laboratório 2}}
\author{{Afonso Cruz (102815)} \and {Eduardo Batista (99931)} \and {Henrique Póvoa (102109)}}
\date{Grupo {nº\textbf{1}} \quad --- \quad 13/11/2026}

\begin{document}
\maketitle

\begin{abstract}
\begin{caixaresumo}
Transpõe-se o swing do Laboratório~1 para um braço real, o UFactory Lite~6
(6~juntas rotacionais). A cinemática direta é obtida pelo método de
Denavit--Hartenberg modificado e validada contra uma segunda tabela do
fabricante e contra o próprio controlador. A ponta da pinça faz de cabeça do
taco: percorre um arco de \SI{150}{\milli\meter} que passa numa bola fixa, com
as três fases e as leis de movimento do Lab~1 e a pinça a inclinar-se como um
taco.

A trajetória usa no máximo \SI{44}{\percent} da velocidade cartesiana e
\SI{53}{\percent} da velocidade das juntas admitidas pelo fabricante, e o
modelo 3D do robô mostra folgas de pelo menos \SI{25}{\milli\meter} entre
elos e \SI{75}{\milli\meter} à mesa. No robô, a metade da velocidade, o
flange segue a referência com \SI{0.5}{\milli\meter} de erro médio. A
cinemática inversa do controlador é exata face ao seu modelo interno
(\SI{2}{\micro\meter}); composta com a nossa cinemática direta dá
\SI{1.4}{\milli\meter}, diferença que vem dos parâmetros geométricos e não da
inversa.
\end{caixaresumo}
\end{abstract}


% =============================================================================
\section{Tarefa 1 --- Cinemática direta}
% =============================================================================

O Lite~6 é um manipulador série com seis juntas rotacionais. As três primeiras
(base, ombro e cotovelo) posicionam o pulso; as três últimas têm eixos que se
cruzam num ponto e formam um pulso esférico, que orienta a ferramenta. A
cinemática direta~\cite{sequeira2024,siciliano2009} dá a pose do flange (posição $p$ e orientação $R$) em função
dos seis ângulos $q = (q_1,\dots,q_6)$.

\subsection{Método e atribuição de referenciais}

Usou-se a convenção de Denavit--Hartenberg \emph{modificada}~\cite{craig2005},
que é a do documento do fabricante~\cite{ufactory_dh} e a que o controlador
usa. A cada elo $i$ fixa-se um referencial $\{i\}$ com as seguintes regras:

\begin{itemize}
  \item o eixo $z_i$ coincide com o eixo de rotação da junta $i$;
  \item o eixo $x_i$ fica sobre a normal comum a $z_i$ e $z_{i+1}$ (se os
        eixos se cruzam, é perpendicular aos dois);
  \item $y_i$ completa o referencial direto;
  \item $\{0\}$ é a base, fixa à mesa, com $z_0$ vertical para cima e a
        origem no tampo, sobre o eixo da junta~1.
\end{itemize}

Os referenciais que resultam estão na Figura~\ref{fig:frames} e na
Tabela~\ref{tab:frames}. Dois pares têm a mesma origem: $\{1\}$ e $\{2\}$ no
ombro, porque os eixos das juntas 1 e 2 se cruzam, e $\{4\}$ e $\{5\}$ no
centro do pulso.

\begin{figure}[H]
\centering
\fig{Photos/referenciais.png}{0.62\textwidth}
\caption{Referenciais D-H modificados com todas as juntas a zero, em vista
lateral (plano $x_0z_0$). Azul: $z_i$, eixo da junta $i$; vermelho: $x_i$.
$\{2\}$ e $\{5\}$ estão desenhados ao lado da sua origem (tracejado), que é a
mesma de $\{1\}$ e de $\{4\}$. Cotas em mm.}
\label{fig:frames}
\end{figure}

\begin{table}[H]
\centering
\caption{Origem e eixo de rotação de cada referencial (com $q = 0$).}
\label{tab:frames}
\small
\begin{tabularx}{\textwidth}{@{}clLl@{}}
\toprule
\rowcolor{tabcab}
Ref. & Junta & Origem & Direção de $z_i$ em $\{0\}$ \\
\midrule
$\{0\}$ & ---      & tampo da mesa, sobre o eixo da junta 1 & $+z_0$ (vertical) \\
$\{1\}$ & base     & ombro, \SI{243.3}{\milli\meter} acima da mesa & $+z_0$ \\
$\{2\}$ & ombro    & a mesma de $\{1\}$ & $+y_0$ (horizontal) \\
$\{3\}$ & cotovelo & \SI{200}{\milli\meter} acima do ombro, na ponta do braço & $-y_0$ \\
$\{4\}$ & antebraço & centro do pulso: \SI{87}{\milli\meter} à frente e \SI{227.6}{\milli\meter} abaixo do cotovelo & $-z_0$ \\
$\{5\}$ & pulso    & a mesma de $\{4\}$ & $+y_0$ \\
$\{6\}$ & flange   & \SI{61.5}{\milli\meter} abaixo do centro do pulso, na face do flange & $-z_0$ \\
\bottomrule
\end{tabularx}
\end{table}

\subsection{Tabela D-H e transformações}

Nesta convenção, passa-se de $\{i-1\}$ para $\{i\}$ com quatro movimentos
elementares, cada um com um parâmetro:
\begin{equation}
\T{i-1}{i} = R_x(\alpha_{i-1})\,T_x(a_{i-1})\,R_z(\theta_i)\,T_z(d_i),
\qquad \theta_i = q_i + \text{offset}_i ,
\label{eq:dh}
\end{equation}
\begin{equation}
\T{i-1}{i} =
\begin{bmatrix}
\cos\theta_i & -\sin\theta_i & 0 & a_{i-1}\\
\sin\theta_i\cos\alpha_{i-1} & \cos\theta_i\cos\alpha_{i-1} & -\sin\alpha_{i-1} & -d_i\sin\alpha_{i-1}\\
\sin\theta_i\sin\alpha_{i-1} & \cos\theta_i\sin\alpha_{i-1} & \cos\alpha_{i-1} & d_i\cos\alpha_{i-1}\\
0 & 0 & 0 & 1
\end{bmatrix}.
\label{eq:elo}
\end{equation}
Os parâmetros do Lite~6 estão na Tabela~\ref{tab:dh}. A pose do flange é o
produto das seis matrizes:
\begin{equation}
\T{0}{6}(q) = \T{0}{1}\,\T{1}{2}\,\T{2}{3}\,\T{3}{4}\,\T{4}{5}\,\T{5}{6}
= \begin{bmatrix} R(q) & p(q)\\ 0 & 1\end{bmatrix}.
\label{eq:fk}
\end{equation}

\begin{table}[H]
\centering
\caption{Parâmetros D-H modificados do Lite~6~\cite{ufactory_dh}.}
\label{tab:dh}
\small
\begin{tabular}{ccccc}
\toprule
\rowcolor{tabcab}
Junta $i$ & offset de $\theta_i$ & $d_i$ [mm] & $\alpha_{i-1}$ & $a_{i-1}$ [mm] \\
\midrule
1 & \ang{0}   & 243,3 & \ang{0}   & 0 \\
2 & \ang{-90} & 0     & \ang{-90} & 0 \\
3 & \ang{-90} & 0     & \ang{180} & 200 \\
4 & \ang{0}   & 227,6 & \ang{90}  & 87 \\
5 & \ang{0}   & 0     & \ang{90}  & 0 \\
6 & \ang{0}   & 61,5  & \ang{-90} & 0 \\
\bottomrule
\end{tabular}
\end{table}

Como o pulso é esférico, a posição separa-se da orientação. O centro do pulso
(origem de $\{4\}$ e $\{5\}$) só depende das três primeiras juntas:
\begin{equation}
p_w = \begin{bmatrix} r\cos q_1\\ r\sin q_1\\ h \end{bmatrix},\qquad
\begin{aligned}
r &= 200\sin q_2 + 87\cos(q_3-q_2) + 227{,}6\,\sin(q_3-q_2),\\
h &= 243{,}3 + 200\cos q_2 + 87\sin(q_3-q_2) - 227{,}6\,\cos(q_3-q_2),
\end{aligned}
\label{eq:pulso}
\end{equation}
e o flange fica \SI{61.5}{\milli\meter} à frente, ao longo do eixo da
ferramenta: $p = p_w + 61{,}5\,R\,\hat e_z$. A orientação exprime-se em
\emph{roll}, \emph{pitch} e \emph{yaw} com
$R = R_z(\text{yaw})\,R_y(\text{pitch})\,R_x(\text{roll})$, a convenção do
SDK~\cite{xarm_sdk}. A implementação está em \texttt{src/kinematics.py}
(funções \texttt{link\_transform} e \texttt{fk}).

A pinça montada (Gripper Lite) não tem desvio nem rotação em relação ao
flange: a sua ponta está a $L = \SI{83.6}{\milli\meter}$ ao longo de $z_6$,
logo $\T{0}{\text{ponta}} = \T{0}{6}\,T_z(L)$.

\subsection{Verificação}

\begin{itemize}[itemsep=3pt]
  \item \textbf{Pose de referência.} Com $q=0$, a equação~\eqref{eq:pulso} dá
        $p_w = (87;\,0;\,215{,}7)$ e $p = (87;\,0;\,154{,}2)$~mm, com
        \emph{roll} de \ang{180} (ferramenta a apontar para baixo), como na
        Figura~\ref{fig:frames}.
  \item \textbf{Segunda tabela do fabricante.} O mesmo documento dá os
        parâmetros na convenção D-H clássica, com outros referenciais.
        Implementada à parte e comparada em \num{10000} configurações
        aleatórias, a maior diferença de posição é
        \SI{2e-13}{\milli\meter}, ou seja, erro de arredondamento.
  \item \textbf{Controlador do robô.} Nos 60 pontos da Tarefa~4, a nossa
        cinemática difere da do controlador em média
        \SI{1.4}{\milli\meter} (máximo \SI{2.0}{\milli\meter}) e
        \ang{0.03} na orientação. A origem desta diferença discute-se na
        Secção~\ref{sec:ik}.
\end{itemize}


% =============================================================================
\section{Tarefa 2 --- Trajetória de swing}
\label{sec:traj}
% =============================================================================

\subsection{Analogia com o Laboratório 1}

No Lab~1 o swing era um pêndulo duplo~\cite{cochran1968}: o braço roda em
torno do ombro e o taco roda em torno do pulso, que fica armado no topo e só
solta perto do impacto. Aqui a ponta da pinça faz de cabeça do taco e o
movimento descreve-se com os mesmos dois ângulos:

\begin{itemize}
  \item $\varphi(t)$, o \textbf{braço}: posição da ponta da pinça num arco
        vertical de raio $R_a = \SI{150}{\milli\meter}$ centrado num
        ``ombro'' virtual por cima da bola; $\varphi = 0$ na bola;
  \item $\psi(t)$, o \textbf{pulso}: inclinação adicional da pinça em relação
        ao arco, armada no topo e solta tarde.
\end{itemize}

A bola está fixa em $b = (190;\,0;\,80)$~mm no referencial da base, à frente
do robô, e o swing faz-se para o lado, na direção $\hat d = \hat y_0$. As
dimensões foram escolhidas para o robô: com o raio de \SI{150}{\milli\meter}
e a bola a \SI{190}{\milli\meter}, todo o arco fica dentro do espaço de
trabalho com a pinça montada.

\begin{destaque}
\textbf{O que o enunciado pede} são as funções do tempo da posição $p(t)$ e
da orientação $R(t)$ da pinça. São as equações~\eqref{eq:pos}
e~\eqref{eq:ori}, com os perfis $\varphi(t)$ e $\psi(t)$ da
Secção~\ref{sec:perfis}.
\end{destaque}

\subsection{Perfis temporais}
\label{sec:perfis}

As três fases e as leis de movimento são as do Lab~1
(Tabela~\ref{tab:fases}), com
$s(\tau) = 10\tau^3 - 15\tau^4 + 6\tau^5$ o polinómio de grau 5 que vai de 0
a 1 com velocidade e aceleração nulas nos extremos.

\begin{table}[H]
\centering
\caption{Fases do swing. $\varphi_{top} = \ang{-90}$,
$\varphi_{end} = \ang{50}$, $\psi_{top} = \ang{30}$; $u$ é o tempo contado
desde o início de cada fase.}
\label{tab:fases}
\small
\begin{tabularx}{\textwidth}{@{}lcLL@{}}
\toprule
\rowcolor{tabcab}
Fase & Duração & Braço $\varphi$ & Pulso $\psi$ \\
\midrule
Backswing & $T_b = \SI{2.0}{\second}$ &
  $\varphi_{top}\,s(u/T_b)$: parte e chega parado &
  $\psi_{top}\,s(u/T_b)$: arma \\
Downswing & $T_d = \SI{2.2}{\second}$ &
  $\varphi_{top}\bigl(1-(u/T_d)^2\bigr)$: aceleração constante, velocidade máxima na bola &
  armado até $0{,}6\,T_d$; depois solta com $s$ e chega a 0 no impacto \\
Follow-through & $T_f = \SI{1.22}{\second}$ &
  $\omega_{imp}\,u - \tfrac12 (\omega_{imp}/T_f)\,u^2$: desaceleração constante até parar em $\varphi_{end}$ &
  0 \\
\bottomrule
\end{tabularx}
\end{table}

A velocidade angular no impacto é $\omega_{imp} = -2\varphi_{top}/T_d =
\SI{1.43}{\radian\per\second}$, que dá \SI{214}{\milli\meter\per\second} na
ponta da pinça, e $T_f = 2\varphi_{end}/\omega_{imp}$. O swing completo dura
\SI{5.42}{\second} e é amostrado a \SI{100}{\hertz} (543 pontos). O
downswing com aceleração constante é o que dá mais velocidade na bola para
uma dada aceleração máxima. A velocidade mais alta de toda a trajetória
(\SI{221}{\milli\meter\per\second}) ocorre a meio do backswing, onde o
polinómio tem o seu máximo; como os tempos foram escolhidos para respeitar os
limites do robô e não para imitar um jogador, o backswing não é mais lento do
que o downswing como num swing real.

\subsection{Posição e orientação da pinça}

\paragraph{Posição.} Com o centro do arco em
$c = b + R_a\,\hat z_0 = (190;\,0;\,230)$~mm, a ponta da pinça é
\begin{equation}
p(t) = c + R_a\bigl(\sin\varphi(t)\,\hat d - \cos\varphi(t)\,\hat z_0\bigr)
     = \bigl(190;\ 150\sin\varphi;\ 230 - 150\cos\varphi\bigr)\ \text{mm}.
\label{eq:pos}
\end{equation}
No topo do backswing está em $(190;\,-150;\,230)$ e no fim do swing em
$(190;\,115;\,134)$~mm.

\paragraph{Orientação.} Na bola, a pinça aponta para baixo e está inclinada
\ang{45} para o lado da base:
$R_0 = R(\text{roll}=\ang{180},\,\text{pitch}=\ang{45},\,\text{yaw}=0)$. O
seu eixo é a terceira coluna de $R_0$, $z_t = (-0{,}707;\,0;\,-0{,}707)$.
Durante o swing roda-se $R_0$ em torno de um eixo fixo $\hat n$, perpendicular
ao eixo da pinça e à direção do swing:
\begin{equation}
R(t) = \operatorname{Rot}\bigl(\hat n,\ \varphi(t)+\psi(t)\bigr)\,R_0,
\qquad
\hat n = \frac{z_t\times\hat d}{\lVert z_t\times\hat d\rVert}
       = (0{,}707;\ 0;\ -0{,}707),
\label{eq:ori}
\end{equation}
com $\operatorname{Rot}(\hat n,\theta) = I + \sin\theta\,K +
(1-\cos\theta)K^2$ (fórmula de Rodrigues, $K$ a matriz do produto externo por
$\hat n$).

Duas escolhas merecem justificação:

\begin{itemize}[itemsep=3pt]
  \item \textbf{Inclinação de \ang{45} na bola.} Com a pinça a pique, a
        junta~5 ficaria em \ang{0}: os eixos das juntas 4 e 6 alinham-se
        (singularidade do pulso) e as suas velocidades disparam, cerca de
        \SI{1900}{\degree\per\second} no primeiro teste. Com \ang{45}, a
        junta~5 fica entre \ang{61} e \ang{80} em todo o swing.
  \item \textbf{Eixo $\hat n$ perpendicular ao eixo da pinça.} A velocidade
        angular é sempre paralela a $\hat n$, logo não tem componente ao longo
        de $z_t$: a pinça inclina-se para o lado, como um taco, sem girar
        sobre si própria. Na primeira versão a rotação era em torno da normal
        ao plano do swing, e \SI{71}{\percent} dela ($\cos\ang{45}$) era a
        pinça a rodar sobre o seu eixo: a junta~6 fazia \ang{124} e a base
        \ang{73}. Agora fazem \ang{39} e \ang{13}, e o flange anda a
        \SI{170}{\milli\meter\per\second} em vez de 298.
\end{itemize}

\paragraph{Flange.} O controlador do laboratório tem o desvio de ferramenta a
zero, portanto comanda o flange e não a ponta. A pose enviada é
\begin{equation}
p_f(t) = p(t) - L\,R(t)\,\hat e_z , \qquad L = \SI{83.6}{\milli\meter},
\end{equation}
com a mesma orientação. Na bola, $p_f = (249{,}1;\,0;\,139{,}1)$~mm. O código
lê o desvio configurado no controlador e converte as poses para o ponto que
ele comanda, pelo que a ponta passa na bola com qualquer configuração.

A Figura~\ref{fig:cart} mostra $p(t)$, a orientação e o módulo da velocidade;
a Figura~\ref{fig:3d} mostra o caminho no espaço.

\begin{figure}[H]
\centering
\fig{Photos/traj_cartesiana.png}{0.74\textwidth}
\caption{Trajetória da ponta da pinça: posição, orientação
(\emph{roll}/\emph{pitch}/\emph{yaw}) e módulo da velocidade da ponta e do
flange, com o limite do fabricante.}
\label{fig:cart}
\end{figure}

\begin{figure}[H]
\centering
\fig{Photos/traj_3d.png}{0.55\textwidth}
\caption{Caminho da ponta da pinça e do flange no referencial da base. Os
segmentos cinzentos são a pinça em instantes sucessivos; os círculos marcam o
cilindro da base.}
\label{fig:3d}
\end{figure}

\subsection{Viabilidade}
\label{sec:viab}

Para verificar as juntas é preciso $q(t)$. Antes de ir ao robô, obteve-se com
uma cinemática inversa numérica própria (mínimos quadrados
amortecidos~\cite{buss2004} sobre o erro de posição e de orientação), ponto a
ponto e partindo sempre da solução anterior, para as juntas variarem de forma
contínua. Esta inversa serve só para a verificação; a do robô é avaliada na
Tarefa~4. As velocidades calculam-se por diferenças centrais. A
Tabela~\ref{tab:viab} resume as três condições do enunciado, com uma margem de
segurança de \SI{20}{\percent} sobre os limites do
fabricante~\cite{ufactory_lite6}; a Tabela~\ref{tab:juntas} e a
Figura~\ref{fig:juntas} detalham as juntas.

\begin{table}[H]
\centering
\caption{Verificação de viabilidade (\texttt{check\_feasibility} em
\texttt{src/trajectory.py}).}
\label{tab:viab}
\small
\begin{tabularx}{\textwidth}{@{}lLccc@{}}
\toprule
\rowcolor{tabcab}
Condição & O que se mede & Limite & Aceite até & Trajetória \\
\midrule
\rowcolor{tabgrupo}\multicolumn{5}{@{}l}{\emph{(i) Velocidade}} \\
Ponta da pinça & $\max\lVert\dot p\rVert$ & \SI{500}{\milli\meter\per\second} & 400 & \textbf{221} \\
Flange         & $\max\lVert\dot p_f\rVert$ & \SI{500}{\milli\meter\per\second} & 400 & \textbf{170} \\
Juntas         & $\max_i\lvert\dot q_i\rvert$ & \SI{180}{\degree\per\second} & 144 & \textbf{95} (junta 4) \\
\midrule
\rowcolor{tabgrupo}\multicolumn{5}{@{}l}{\emph{(ii) Alcance}} \\
Ombro ao centro do pulso & $\max\lVert p_w - p_{\{2\}}\rVert$ & \SI{443.7}{\milli\meter} & 421,5 & \textbf{299} \\
\midrule
\rowcolor{tabgrupo}\multicolumn{5}{@{}l}{\emph{(iii) Limites das juntas}} \\
Folga ao limite mais próximo & $\min_i$ (Tabela~\ref{tab:juntas}) & --- & $\geq\ang{5}$ & \textbf{\ang{44}} (junta 5) \\
\bottomrule
\end{tabularx}
\end{table}

O alcance mede-se do eixo do ombro ao centro do pulso, porque é essa
distância que não pode passar o comprimento do braço esticado,
$200 + \sqrt{87^2 + 227{,}6^2} = \SI{443.7}{\milli\meter}$; perto desse valor
o cotovelo entra em singularidade, e por isso só se aceita até
\SI{95}{\percent}. A trajetória usa \SI{67}{\percent}.

\begin{table}[H]
\centering
\caption{Limites das juntas~\cite{ufactory_lite6} e gama usada no swing.}
\label{tab:juntas}
\small
\begin{tabular}{clccc}
\toprule
\rowcolor{tabcab}
Junta & Papel & Limites & Gama usada & Amplitude \\
\midrule
1 & base      & $\pm\ang{360}$            & \ang{-12} a \ang{1}  & \ang{13} \\
2 & ombro     & $\pm\ang{150}$            & \ang{16} a \ang{48}  & \ang{32} \\
3 & cotovelo  & \ang{-3.5} a \ang{300}    & \ang{45} a \ang{63}  & \ang{18} \\
4 & antebraço & $\pm\ang{360}$            & \ang{-66} a \ang{55} & \ang{121} \\
5 & pulso     & $\pm\ang{124}$            & \ang{61} a \ang{80}  & \ang{19} \\
6 & flange    & $\pm\ang{360}$            & \ang{-26} a \ang{14} & \ang{39} \\
\bottomrule
\end{tabular}
\end{table}

\begin{figure}[H]
\centering
\fig{Photos/traj_juntas.png}{0.72\textwidth}
\caption{Ângulos e velocidades das seis juntas ao longo do swing. A tracejado,
o limite de velocidade do fabricante ($\pm\SI{180}{\degree\per\second}$). A
junta~4 (rotação do antebraço) é a que faz a inclinação lateral da pinça.}
\label{fig:juntas}
\end{figure}

\subsection{Ausência de contacto perigoso}
\label{sec:contacto}

A verificação anterior inclui duas condições geométricas simples: a ponta, o
flange e as origens de todos os referenciais ficam pelo menos
\SI{30}{\milli\meter} acima da mesa (o mínimo é \SI{80}{\milli\meter}, a
própria bola) e a mais de \SI{80}{\milli\meter} do eixo da base (o mínimo é
\SI{190}{\milli\meter}). Estas condições só olham para pontos; não veem o
volume dos elos nem da pinça.

Para o robô inteiro, reproduziu-se a trajetória das juntas no
MuJoCo~\cite{todorov2012}, com o modelo oficial do Lite~6 com
pinça~\cite{menagerie2022}, e mediu-se em cada instante a distância entre as
malhas de colisão de todos os pares de corpos que podem tocar-se e entre o
robô e o plano da mesa (\texttt{src/mujoco\_preview.py}). As malhas de
colisão são invólucros convexos dos elos, logo as distâncias são
conservadoras. A Figura~\ref{fig:seq} mostra o movimento e a
Figura~\ref{fig:folgas} as folgas.

\begin{figure}[H]
\centering
\fig{Photos/mj_sequencia_down_the_line.png}{\textwidth}
\caption{O swing no modelo 3D do robô, visto de trás da linha de jogo. A
vermelho, o caminho da ponta da pinça; a esfera é a bola.}
\label{fig:seq}
\end{figure}

\begin{figure}[H]
\centering
\fig{Photos/mj_folgas.png}{0.72\textwidth}
\caption{Folga mínima ao longo do swing: entre o robô (com pinça) e a mesa, e
entre elos não adjacentes. As linhas verticais marcam o topo do backswing e o
impacto.}
\label{fig:folgas}
\end{figure}

\begin{itemize}[itemsep=3pt]
  \item \textbf{Mesa.} O ponto mais baixo do robô é a pinça, no endereço e
        no impacto, a \SI{75}{\milli\meter} da mesa.
  \item \textbf{Entre elos.} A menor distância é de \SI{25}{\milli\meter},
        entre o corpo da pinça e o elo~4, e é praticamente constante: é a
        distância construtiva entre os dois, que nesta trajetória quase não
        muda. O par seguinte (elos 2 e 4) nunca fica a menos de
        \SI{34}{\milli\meter}.
  \item \textbf{Modelo.} O modelo 3D e a nossa cinemática coincidem a menos de
        \SI{1}{\milli\meter} no flange e \SI{2}{\milli\meter} na ponta da
        pinça, diferença desprezável face às folgas.
\end{itemize}

A bola é um ponto fixo no espaço de trabalho, como o enunciado permite; não
há nenhum objeto na mesa, que está livre à volta do robô. A ponta da pinça
passa nesse ponto a \SI{214}{\milli\meter\per\second}, a
\SI{80}{\milli\meter} do tampo.


% =============================================================================
\section{Tarefa 3 --- Execução no robô}
\label{sec:exec}
% =============================================================================

\subsection{Script}

O script \texttt{src/run\_robot.py} usa só funções do
SDK~\cite{xarm_sdk}. A sequência é:

\begin{enumerate}[itemsep=2pt]
  \item gera a trajetória e corre a verificação da Secção~\ref{sec:viab}; se
        alguma condição falhar, nada é enviado;
  \item liga ao robô, lê o desvio de ferramenta do controlador e converte as
        poses para o ponto que ele comanda; pede confirmação ao operador;
  \item aproxima-se da pose inicial junta a junta, a
        \SI{20}{\degree\per\second} (\texttt{set\_servo\_angle}), para não
        atravessar zonas imprevistas a partir de uma pose qualquer;
  \item executa o swing por um de dois métodos.
\end{enumerate}

\paragraph{Método \texttt{servo}.} Em modo~1, envia uma pose a cada
\SI{10}{\milli\second} com \texttt{set\_servo\_cartesian}
(Listagem~\ref{lst:servo}). É o controlador que resolve a cinemática inversa
em cada ponto, e o perfil temporal é o nosso: é o método que preserva a
aceleração constante do downswing. Em cada ciclo lê-se a pose com
\texttt{get\_position} e grava-se ao lado da pedida.

\paragraph{Método \texttt{pontos}.} Em modo~0, envia um ponto em cada dez com
\texttt{set\_position} e raio de concordância de \SI{5}{\milli\meter}; o
controlador interpola entre eles. É mais simples, mas o perfil de velocidade
passa a ser o do controlador.

O parâmetro \texttt{-{}-lento N} estica o tempo $N$ vezes, mantendo o
caminho; permite ensaiar devagar antes de passar à velocidade planeada.

\begin{lstlisting}[language=Python, caption={Ciclo do método \texttt{servo}
(simplificado de \texttt{run\_servo} em \texttt{src/run\_robot.py}).},
label={lst:servo}, captionpos=b]
arm.set_mode(1); arm.set_state(0)            # modo servo
t0 = time.perf_counter()
for k, pose in enumerate(poses):             # [x, y, z, roll, pitch, yaw]
    if arm.set_servo_cartesian(pose) != 0:   # erro do controlador: parar
        break
    _, medida = arm.get_position()
    registo.append([time.perf_counter() - t0, *pose, *medida])
    while time.perf_counter() - t0 < (k + 1) * dt:
        pass                                 # espera até ao instante seguinte
arm.set_mode(0); arm.set_state(0)
\end{lstlisting}

\subsection{Resultados}

Os ensaios da versão final fizeram-se a 6/11/2026 (controlador com
\emph{firmware} 2.7.1). A Tabela~\ref{tab:exec} resume o seguimento e a
Figura~\ref{fig:exec} compara a posição pedida com a medida. O erro é a
distância entre o flange medido e o pedido no mesmo instante; o atraso é o
deslocamento temporal que minimiza esse erro.

\begin{table}[H]
\centering
\caption{Seguimento da trajetória em modo \texttt{servo} (posição do flange).
Cada ensaio foi feito duas vezes.}
\label{tab:exec}
\small
\begin{tabular}{lcccccc}
\toprule
\rowcolor{tabcab}
Ensaio & Duração & Pontos & Erro médio & Erro máx. & Atraso & Erro sem o atraso \\
\midrule
\texttt{-{}-lento 2}, 1.º & \SI{10.8}{\second} & 1085/1085 & \SI{0.54}{\milli\meter} & \SI{2.5}{\milli\meter} & \SI{10}{\milli\second} & \SI{0.17}{\milli\meter} \\
\texttt{-{}-lento 2}, 2.º & \SI{10.9}{\second} & 1085/1085 & \SI{0.57}{\milli\meter} & \SI{3.2}{\milli\meter} & \SI{10}{\milli\second} & \SI{0.27}{\milli\meter} \\
\texttt{-{}-lento 1}, 1.º & \SI{5.4}{\second}  & 543/543   & \SI{8.4}{\milli\meter}  & \SI{28.0}{\milli\meter} & \SI{58}{\milli\second} & \SI{7.2}{\milli\meter} \\
\texttt{-{}-lento 1}, 2.º & \SI{5.4}{\second}  & 543/543   & \SI{8.4}{\milli\meter}  & \SI{28.2}{\milli\meter} & \SI{55}{\milli\second} & \SI{7.2}{\milli\meter} \\
\bottomrule
\end{tabular}
\end{table}

\begin{figure}[H]
\centering
\fig{Photos/run_servo.png}{\textwidth}
\caption{Posição do flange pedida (tracejado) e medida, e erro de posição, a
metade da velocidade (esquerda) e à velocidade planeada (direita). As linhas
verticais separam as três fases.}
\label{fig:exec}
\end{figure}

\begin{itemize}[itemsep=3pt]
  \item \textbf{A metade da velocidade o robô segue a trajetória.} O erro
        médio é de \SI{0.5}{\milli\meter}, da ordem da repetibilidade do robô
        ($\pm\SI{0.5}{\milli\meter}$~\cite{ufactory_lite6}), e quase todo é
        um atraso de um ciclo (\SI{10}{\milli\second}). No impacto o flange
        está a \SI{1.3}{\milli\meter} da referência.
  \item \textbf{À velocidade planeada o robô atrasa-se onde a velocidade
        varia depressa.} O erro concentra-se no backswing (média
        \SI{11}{\milli\meter}) e no follow-through (\SI{16}{\milli\meter});
        no downswing, onde a aceleração é constante, é de
        \SI{1.4}{\milli\meter}. Não é um atraso puro: descontar
        \SI{58}{\milli\second} só baixa o erro de 8,4 para
        \SI{7.2}{\milli\meter}. A velocidade máxima medida
        (\SI{173}{\milli\meter\per\second}) iguala a pedida, por isso não é
        saturação de velocidade; o comportamento é compatível com um limite
        de aceleração ou uma filtragem do controlador em modo servo.
        \todo{não confirmámos qual; ver se o Studio mostra limites de
        aceleração.}
  \item \textbf{O resultado é reprodutível.} Dois ensaios iguais dão o mesmo
        erro (Tabela~\ref{tab:exec}).
  \item \textbf{O método \texttt{pontos}} executou o swing completo com
        \texttt{-{}-lento 3}, sem erros. Não regista a pose medida.
\end{itemize}

\subsection{Ocorrências e correções}

Na primeira sessão (5/11) a trajetória era diferente em dois pontos: o
follow-through ia até $\varphi_{end} = \ang{70}$ e a pinça rodava em torno da
normal ao plano. Três ocorrências levaram à versão final:

\begin{itemize}[itemsep=3pt]
  \item \textbf{Erro C31 (deteção de colisão).} Com $\varphi_{end}=\ang{70}$
        o robô parou sempre no mesmo ponto, a $\varphi \approx \ang{69}$,
        quase parado e com o braço esticado para o lado (junta~5 a
        \ang{110}). Sendo num ponto fixo e a baixa velocidade, não é um
        efeito dinâmico; as hipóteses são o cabo
        da pinça a esticar ou a massa da pinça não estar configurada no
        controlador. Encurtar o follow-through para \ang{50} eliminou o erro
        em todos os ensaios seguintes. \todo{correr \texttt{robot\_info.py}
        no laboratório para ler a carga configurada (\texttt{tcp\_load}).}
  \item \textbf{Erro C24 (velocidade excedida)} no método \texttt{pontos} à
        velocidade planeada: a concordância entre pontos pede às juntas mais
        do que o plano. O script passou a recusar este método com
        $N < 2$.
  \item \textbf{Seguimento fraco com a primeira orientação.} O flange tinha
        de andar a \SI{298}{\milli\meter\per\second} e o erro médio era de
        \SI{2.0}{\milli\meter} a metade da velocidade e
        \SI{18.9}{\milli\meter} à velocidade planeada. Com a pinça a
        inclinar-se sem girar (Secção~\ref{sec:traj}), o erro desceu para
        0,5 e \SI{8.4}{\milli\meter}.
\end{itemize}

A verificação de viabilidade também foi corrigida depois desta sessão: só
olhava para a velocidade da ponta da pinça, quando o ponto que o controlador
comanda é o flange.


% =============================================================================
\section{Tarefa 4 --- Precisão da cinemática inversa do robô}
\label{sec:ik}
% =============================================================================

\subsection{Método}

Escolheram-se 60 pontos igualmente espaçados no tempo ao longo de todo o swing
(\texttt{src/ik\_compare.py}). Para cada pose pedida $(p_k, R_k)$:

\begin{enumerate}[itemsep=2pt]
  \item $q_k$ = \texttt{get\_inverse\_kinematics}$(p_k, R_k)$, a inversa do
        controlador;
  \item $(\hat p_k, \hat R_k) = \T{0}{6}(q_k)$, a nossa cinemática direta da
        Tarefa~1;
  \item erro de posição $\lVert\hat p_k - p_k\rVert$ e erro de orientação, o
        ângulo da rotação $\hat R_k R_k^{T}$.
\end{enumerate}

Para separar o erro da inversa do erro do nosso modelo, calculou-se também a
pose com a cinemática direta do próprio controlador
(\texttt{get\_forward\_kinematics}). O robô não se move neste ensaio.

\subsection{Resultados}

\begin{table}[H]
\centering
\caption{Pose recuperada a partir das juntas da inversa do robô, comparada com
a pose pedida (60 pontos).}
\label{tab:ik}
\small
\begin{tabular}{lcccc}
\toprule
\rowcolor{tabcab}
 & \multicolumn{2}{c}{Posição} & \multicolumn{2}{c}{Orientação} \\
\rowcolor{tabcab}
Cinemática direta usada & média & máx. & média & máx. \\
\midrule
Do controlador & \SI{0.1}{\micro\meter} & \SI{1.9}{\micro\meter} & --- & \ang{0.0004} \\
A nossa (Tabela~\ref{tab:dh}) & \SI{1.36}{\milli\meter} & \SI{2.02}{\milli\meter} & \ang{0.014} & \ang{0.030} \\
\bottomrule
\end{tabular}
\end{table}

\begin{figure}[H]
\centering
\fig{Photos/ik_erros.png}{0.8\textwidth}
\caption{Erro da pose recuperada ao longo do swing. Em cima: erro de posição
com a nossa cinemática direta e com a do controlador. Ao meio: componentes da
diferença com a nossa cinemática. Em baixo: erro de orientação. Os círculos
marcam os pontos em que a inversa devolveu o pulso virado.}
\label{fig:ik}
\end{figure}

\subsection{Discussão}

\begin{itemize}[itemsep=3pt]
  \item \textbf{A inversa do robô é exata face ao seu modelo.} Resolveu os 60
        pontos, e a cinemática direta do controlador devolve a pose pedida
        com menos de \SI{2}{\micro\meter} e \ang{0.0004}, ou seja, ao
        nível do arredondamento.
  \item \textbf{A diferença de \SI{1.4}{\milli\meter} está no modelo, não na
        inversa.} É, ponto a ponto, a diferença entre a nossa cinemática
        direta e a do controlador. Tem um desvio sistemático de
        $(-0{,}2;\,+0{,}6;\,-0{,}9)$~mm e a componente vertical quase não
        varia ao longo do swing. O erro de orientação é muito pequeno: um
        desvio de \ang{0.03} num braço de \SI{300}{\milli\meter} desloca a
        ponta \SI{0.16}{\milli\meter}, dez vezes menos do que a diferença
        medida. A diferença vem, por isso, sobretudo de comprimentos e não
        de ângulos.
  \item \textbf{A diferença é compatível com parâmetros geométricos
        ligeiramente diferentes dos nominais.} Ajustando pequenas correções
        aos parâmetros da Tabela~\ref{tab:dh} para reproduzir o controlador
        nestes 60 pontos, a diferença desce de 1,36 para
        \SI{0.06}{\milli\meter}. Os mesmos parâmetros, aplicados a dados que
        não entraram no ajuste (juntas e pose medidas
        \ifsegundoswing no swing da Secção~\ref{sec:v3}, \else num segundo
        swing executado no robô, desenhado no espaço das juntas, \fi noutra
        zona do espaço de trabalho), baixam a
        diferença de 1,6 para \SI{0.9}{\milli\meter}
        (\texttt{src/ajuste\_dh.py}). O controlador usa, portanto, um modelo
        cinemático que não é exatamente o da tabela nominal. Os valores
        ajustados não se podem tomar como a calibração do robô: com poses de
        uma só trajetória, vários parâmetros compensam-se entre si.
        \todo{ler os parâmetros D-H do controlador com
        \texttt{robot\_info.py} para fechar este ponto.}
  \item \textbf{A inversa não garante continuidade das juntas.} Em 19 dos 60
        pontos devolveu a solução de pulso virado
        ($q_4 \pm \ang{180}$, $-q_5$, $q_6 \pm \ang{180}$), que dá a mesma
        pose com outros ângulos. Nos restantes 41, e nos 19 depois de
        convertidos, coincide com a nossa inversa numérica a menos de
        \ang{0.5} em todas as juntas. Na Figura~\ref{fig:ik} vê-se que a
        diferença para o nosso modelo depende do ramo: nos pontos de pulso
        virado a componente $y$ quase se anula.
\end{itemize}

\begin{destaque}
\textbf{Conclusão.} A cinemática inversa incorporada é numericamente exata e
resolveu todos os pontos da trajetória. A precisão com que a ferramenta
atinge uma pose é limitada pelo modelo geométrico e pela repetibilidade do
robô, da ordem de 1 a \SI{2}{\milli\meter}, e não pela inversa. Para executar
uma trajetória com as juntas obtidas ponto a ponto, é preciso escolher o ramo
da solução, porque a inversa pode alternar entre soluções equivalentes.
\end{destaque}


\ifsegundoswing
% =============================================================================
\section{Um segundo swing, desenhado nas juntas}
\label{sec:v3}
% =============================================================================

O swing da Secção~\ref{sec:traj} é pequeno: a base roda \ang{13}. Para um
movimento mais próximo do de um jogador, em que o tronco roda, fez-se um
segundo swing diretamente no espaço das juntas (\texttt{src/swing\_v3.py}).

Moveu-se o robô à mão (modo manual) ao longo do gesto pretendido e
leram-se as juntas em 20 instantes. Dessa demonstração tiraram-se três poses: topo do
backswing, impacto e fim. O caminho entre elas foi refeito com as leis do
Lab~1, em todas as juntas ao mesmo tempo: polinómio de grau~5 no backswing;
aceleração constante no downswing para as juntas que fazem o swing (base,
pulso e rotação da pinça), que chegam ao impacto no mesmo instante e à
velocidade máxima, com o pulso armado durante a primeira metade; e um
polinómio de grau~5 até parar no follow-through. A base roda \ang{125}. A
posição e a orientação da pinça obtêm-se pela cinemática direta da Tarefa~1
(Figura~\ref{fig:v3}), e os tempos foram ajustados para a ponta não passar de
\SI{400}{\milli\meter\per\second}.

\begin{figure}[H]
\centering
\fig{Photos/v3_caminho.png}{0.9\textwidth}
\caption{Segundo swing. Esquerda: caminho da ponta da pinça e a demonstração
feita à mão (a laranja). Direita: velocidade da ponta e do flange à
velocidade planeada; a linha vertical marca o impacto.}
\label{fig:v3}
\end{figure}

\begin{figure}[H]
\centering
\fig{Photos/mjv3demo_sequencia_face_on.png}{\textwidth}
\caption{O segundo swing no modelo 3D do robô.}
\label{fig:v3seq}
\end{figure}

É enviado junta a junta com \texttt{set\_servo\_angle\_j}, sem cinemática
inversa. A metade da velocidade (\SI{20.4}{\second}, 2044 pontos, dois
ensaios) o erro médio das juntas foi de \ang{0.05} (máximo \ang{0.6}) e a
ponta da pinça ficou em média a \SI{1.0}{\milli\meter} da planeada. No modelo
3D a folga mínima é de \SI{17}{\milli\meter}, entre os dedos da pinça e a
mesa logo a seguir ao impacto, e de \SI{21}{\milli\meter} entre elos: não há
contacto, mas a margem à mesa é mais curta do que os \SI{30}{\milli\meter}
exigidos ao primeiro swing. \todo{esta secção só entra se o swing for
repetido no robô com a ponta mais alta (\texttt{-{}-dj2 -6}: 33 mm de folga à
mesa no modelo 3D); atualizar então os números e as figuras.}
\fi


% =============================================================================
\section{Comentário aos resultados}
\label{sec:comentario}
% =============================================================================

\subsection*{\textcolor{bom}{Pontos bons}}

\begin{itemize}[itemsep=3pt]
  \item \textbf{Cinemática direta validada de três formas}: pose de
        referência, segunda tabela do fabricante (\SI{e-13}{\milli\meter}) e
        controlador (\SI{1.4}{\milli\meter}).
  \item \textbf{Trajetória com folga em todos os limites.} Usa
        \SI{44}{\percent} da velocidade cartesiana, \SI{53}{\percent} da
        velocidade das juntas e \SI{67}{\percent} do alcance, e a junta mais
        próxima de um limite fica a \ang{44} dele.
  \item \textbf{Ausência de contacto verificada com o volume real do robô},
        e não só com pontos: \SI{25}{\milli\meter} entre elos e
        \SI{75}{\milli\meter} à mesa.
  \item \textbf{Seguimento de \SI{0.5}{\milli\meter}} a metade da velocidade,
        reprodutível entre ensaios.
  \item \textbf{A escolha do eixo de rotação da pinça} reduziu o erro de
        seguimento quatro vezes e eliminou o movimento desnecessário das
        juntas 1 e 6.
\end{itemize}

\subsection*{\textcolor{mau}{O que não correu tão bem}}

\begin{itemize}[itemsep=3pt]
  \item \textbf{À velocidade planeada o robô não segue o perfil}: erro médio
        de \SI{8}{\milli\meter} e máximo de \SI{28}{\milli\meter}, apesar de
        a trajetória respeitar todos os limites publicados. A verificação de
        viabilidade não inclui acelerações.
  \item \textbf{A causa do erro C31 não ficou identificada}; foi contornada
        encurtando o follow-through.
  \item \textbf{O swing é pequeno e lento} face ao do Lab~1: arco de
        \SI{150}{\milli\meter} e \SI{0.21}{\meter\per\second} no impacto. É
        uma consequência das dimensões do robô e do limite de
        \SI{500}{\milli\meter\per\second}.
  \item \textbf{A diferença de \SI{1.4}{\milli\meter}} para o controlador
        ficou explicada qualitativamente, mas não se identificaram os
        parâmetros que a causam.
\end{itemize}

\subsection*{\textcolor{azul}{O que podíamos melhorar}}

\begin{itemize}[itemsep=3pt]
  \item \textbf{Limitar a aceleração e o \emph{jerk}} no desenho dos perfis,
        para o robô seguir a trajetória à velocidade planeada.
  \item \textbf{Ler a configuração do controlador} (carga da ferramenta,
        parâmetros cinemáticos, sensibilidade de colisão) para fechar as
        duas questões em aberto.
  \item \textbf{Calibrar a cinemática} com poses medidas em todo o espaço de
        trabalho, em vez de uma só trajetória.
  \item \textbf{Desenhar o swing nas juntas}, com a base a fazer de tronco
        \ifsegundoswing como no segundo swing, mantendo as margens de
        segurança do primeiro.\else como num jogador, para um movimento mais
        amplo do que os \ang{13} de rotação da base deste swing.\fi
\end{itemize}

\newpage
% =============================================================================
\section*{Contribuição dos membros}
\todo{rever esta secção; o texto abaixo é uma proposta.}
Todos os membros contribuíram de forma igual para a realização do projeto.
Cada um ficou responsável por uma área, sem deixar de participar nas restantes:
\begin{itemize}
  \item \textbf{Afonso Cruz:} verificação de contacto no modelo 3D e
        redação do relatório.
  \item \textbf{Eduardo Batista:} \todo{área.}
  \item \textbf{Henrique Póvoa:} código de cinemática, trajetória e execução
        no robô; ensaios no laboratório.
\end{itemize}

% =============================================================================
\bibliographystyle{IEEEtran}
\bibliography{refs}

\end{document}
